\documentclass[../main.tex]{subfiles}
\begin{document}
\lstnewenvironment{latexexample}{\lstset{language={[LaTeX]TeX},numbers=none}}{}
\section*{参考：コード記述テンプレート（暫定）}
このページは執筆時にコピーして使う記述例であり、研究結果ではない。
共通設定の \texttt{listings}、\texttt{jvlisting}、\texttt{needspace} を使用する。

\subsection*{インラインコード}
記述例：\verb|\lstinline!cudaDeviceSynchronize()!|。
表示例：\lstinline!cudaDeviceSynchronize()! を呼び出す。
アンダースコアを含む識別子も \lstinline!device_buffer! のように記述できる。
区切り文字はコード中に含まれない文字を選ぶ。

\subsection*{コードブロック}
次の記述を本文へコピーし、言語、説明、ラベル、コードを置き換える。

\begin{latexexample}
\begin{lstlisting}[language=C++,caption={同期処理の例},label={lst:sync-example}]
// ホスト側で完了を待つ
cudaDeviceSynchronize();
\end{lstlisting}
\end{latexexample}

表示例：
\begin{lstlisting}[language=C++,caption={同期処理の例（テンプレート）},label={lst:template-sync}]
// ホスト側で完了を待つ
cudaDeviceSynchronize();
\end{lstlisting}
本文からの参照は \verb|リスト\ref{lst:template-sync}| と記述する。
表示例：リスト\ref{lst:template-sync}。
ラベルは各コードブロックで重複しない名前にする。

\subsection*{短いコードの直前で改ページを調整する}
直後に指定行数分の余白がない場合、\verb|\Needspace| で改ページを促す。
コード全体を囲む命令ではなく、長いコードのページ分割を禁止するものでもない。

\begin{lstlisting}[language={[LaTeX]TeX},numbers=none]
\Needspace{8\baselineskip}
% この直後にコードブロックを置く。
\end{lstlisting}

\Needspace{8\baselineskip}
\begin{lstlisting}[language=C++,caption={短いコードの例（テンプレート）},label={lst:template-short}]
int main() {
    // 実際の処理へ置き換える
    return 0;
}
\end{lstlisting}
\end{document}
